\section*{Hoofdstuk \ref{ch:b3:supramolecular} --- Supramoleculaire chemie}

\begin{solution}{exo:b3:supramolecular:1}
Ion-dipool; $\pi$-stapeling; kation-$\pi$; \omterm{def:b3:supramolecular:interactions}{halogeenbinding} (I naar N); drie waterstofbruggen.
\end{solution}

\begin{solution}{exo:b3:supramolecular:2}
$\Delta G^\circ = -RT\ln K = -8.314 \times 298 \times \ln(\num{1.0e4}) = \qty{-22.8}{kJ/mol}$.
\end{solution}

\begin{solution}{exo:b3:supramolecular:3}
$K_{\ce{K+}}/K_{\ce{Na+}} = 10^{1.7} = 50$; het verschil is $RT\ln 50 = \qty{9.7}{kJ/mol}$ in het voordeel van kalium.
\end{solution}

\begin{solution}{exo:b3:supramolecular:4}
$n/(n + m) = 1/3$: één \omterm{def:b3:supramolecular:supramolecular}{gastheer} voor twee \omterm{def:b3:supramolecular:supramolecular}{gasten}, $\mathrm{HG_2}$.
\end{solution}

\begin{solution}{exo:b3:supramolecular:5}
$H_0 + G_0 + 1/K = \qty{4.667e-3}{mol/L}$; $[\mathrm{HG}] = (4.667 - \sqrt{4.667^2 - 4 \times 2.0 \times 2.0})/2 \times 10^{-3} = \qty{1.13e-3}{mol/L}$: gebonden
fractie 0.566, $\delta = 3.560 + 0.160 \times 0.566 = \qty{3.651}{ppm}$.
\end{solution}

\begin{solution}{exo:b3:supramolecular:6}
$\log K = 18.3 - 8.6 = 9.7$, $K = \num{5e9}$; $\Delta G^\circ = -RT\ln K = \qty{-55}{kJ/mol}$. De zes bindingen Ni--N zijn aan
beide kanten van hetzelfde soort; het aantal vrije moleculen stijgt van
vier naar zeven: de winst aan translatie-entropie drijft de reactie.
\end{solution}

\begin{solution}{exo:b3:supramolecular:7}
$\Delta G^\circ = -RT\ln(\num{2.0e5}) = \qty{-30.2}{kJ/mol}$; $T\Delta S^\circ = \Delta H^\circ - \Delta G^\circ = -40 + 30.2 = \qty{-9.8}{kJ/mol}$;
$\Delta S^\circ = \qty{-33}{J.K^{-1}.mol^{-1}}$: enthalpiegedreven binding, deels terugbetaald in entropie.
\end{solution}

\begin{solution}{exo:b3:supramolecular:8}
Koper(I) bindt twee tweetandige fenantroline-eenheden en houdt ze, omdat
het tetraëdrisch is, loodrecht op elkaar, de ene geregen door de ruimte
waar de ring van de andere zich zal sluiten. Het sluiten van beide ketens
tot ringen vervlecht ze dan. Een grote overmaat cyanide, dat koper(I) veel
sterker bindt, verwijdert het metaal en laat het vrije \omterm{def:b3:supramolecular:interlocked}{catenaan} over.
\end{solution}

\begin{solution}{exo:b3:supramolecular:9}
$k_c = \pi\Delta\nu/\sqrt2 = \pi \times 120/1.414 = \qty{267}{s^{-1}}$. Eyring:
$\Delta G^\ddagger = RT\ln(k_BT/hk_c) = 8.314 \times 300 \times \ln(\num{6.25e12}/267) = \qty{59.6}{kJ/mol}$.
\end{solution}

\begin{solution}{exo:b3:supramolecular:10}
Met een overmaat vast geneesmiddel ligt het vrije geneesmiddel vast op $S_0$. Dan is $[\mathrm{DC}] = KS_0[\mathrm C]$ en
$C_t = [\mathrm C] + [\mathrm{DC}] = [\mathrm C](1 + KS_0)$, dus $[\mathrm{DC}] = KS_0C_t/(1 + KS_0)$ en de totale oplosbaarheid is $S_0 + [\mathrm{DC}]$. Hier is
$KS_0 = 0.050$: $S = 0.10 + 0.050 \times 10/1.050 = \qty{0.58}{mmol/L}$, bijna zes keer de intrinsieke waarde.
\end{solution}

\begin{solution}{exo:b3:supramolecular:11}
Als $1/K \ll H_0, G_0$, is $[\mathrm{HG}] \approx \bigl(H_0 + G_0 - |H_0 - G_0|\bigr)/2 = \min(H_0, G_0)$: de gebonden fractie is $G_0/H_0$ tot één
equivalent, daarna 1. De kromme hangt niet meer af van $K$: elke waarde die groot genoeg is,
geeft binnen de ruis dezelfde scherpe knik, zodat de gegevens alleen tonen dat $K$ groter is dan
een zeker minimum.
\end{solution}

\begin{solution}{exo:b3:supramolecular:12}
$k_{\mathrm{off}} = k_{\mathrm{on}}/K = \qty{1e9}{L.mol^{-1}.s^{-1}}/\qty{1500}{L.mol^{-1}} = \qty{6.7e5}{s^{-1}}$. Het verschil in verschuiving
$\qty{0.160}{ppm}$ is \qty{64}{Hz} bij \qty{400}{MHz}; coalescentie zou $k \approx \pi \times 64/\sqrt2 = \qty{140}{s^{-1}}$ vereisen. De uitwisseling is
duizenden keren sneller: één gemiddeld signaal.
\end{solution}

\begin{solution}{pb:b3:supramolecular:1}
\textbf{1.} Een ring van zes eenheden $\ce{-CH2CH2O-}$; in het complex zit $\ce{K+}$ in het centrum, gebonden door de
zes zuurstofatomen die naar binnen wijzen, met de groepen $\ce{CH2}$ aan de buitenkant.

\textbf{2.} $\ce{K+}$: $-RT\ln 10 \times 6.1 = \qty{-34.8}{kJ/mol}$; $\ce{Na+}$: \qty{-25.1}{kJ/mol}.

\textbf{3.} $10^{6.1 - 4.4} = 50$.

\textbf{4.} Het grotere $\ce{K+}$ (\qty{138}{pm}) bereikt alle zes zuurstofatomen van de ring tegelijk; het
kleinere $\ce{Na+}$ (\qty{102}{pm}) kan dat niet, tenzij de ring plooit, wat spanning kost.

\textbf{5.} Zijn buitenkant is een koolwaterstofoppervlak en zijn lading
is verborgen; $\ce{KCl}$
zou zijn ionen gesolvateerd moeten krijgen, wat tolueen niet kan.

\textbf{6.} Water solvateert $\ce{K+}$ veel beter dan methanol, en het ion moet dat
water kwijtraken om de ring binnen te gaan.

\textbf{7.} $\ce{K+}(\text{aq}) + \ce{MnO4-}(\text{aq}) + \mathrm L(\text{org}) \rightleftharpoons \ce{[KL]+MnO4-}(\text{org})$.

\textbf{8.} $[\ce{K+}]_{\mathrm{aq}}$ blijft \qty{0.10}{mol/L}: $y = \num{1.0e5} \times 0.10 \times (\num{1.0e-3} - y)^2 = \num{1.0e4}(\num{1.0e-3} - y)^2$.

\textbf{9.} Met $u = \num{1.0e-3} - y$: $\num{1.0e4}u^2 + u - \num{1.0e-3} = 0$, $u = \qty{2.70e-4}{mol/L}$, $y = \qty{7.30e-4}{mol/L}$: 73\,\%
geëxtraheerd.

\textbf{10.} $y = \num{1.0e4}(\num{1.0e-3} - y)(\num{1.0e-2} - y)$ geeft $y = \qty{9.89e-4}{mol/L}$: 99\,\%.

\textbf{11.} $y = \num{1.0e4}(\num{1.0e-3} - y)(\num{1.0e-4} - y)$ geeft $y = \qty{9.0e-5}{mol/L}$: 9\,\%, beperkt door de \omterm{def:b3:supramolecular:hosts}{kroonether},
die voor 90\,\% beladen is.

\textbf{12.} Het permanganaation behoudt zijn paarse kleur in het tolueen.
Daar is het niet gesolvateerd, een naakt anion, en reageert het veel
sneller dan in water, waar een hydratatiemantel het afschermt.

\textbf{13.} Het natriumion komt en gaat veel sneller dan het verschil van
de twee resonantiefrequenties: snelle uitwisseling.

\textbf{14.} $\delta = \delta_{\mathrm{vrij}} + (\delta_{\mathrm{geb}} - \delta_{\mathrm{vrij}})[\mathrm{HG}]/H_0$, met $[\mathrm{HG}]$ uit de exacte isotherm, $H_0 = \qty{2.0e-3}{mol/L}$,
$G_0 = (\text{equiv.})\,H_0$, $\delta_{\mathrm{vrij}} = \qty{3.560}{ppm}$.

\textbf{15.} $K = (1.55 \pm 0.08) \times 10^3\ \unit{L.mol^{-1}}$, $\delta_{\mathrm{geb}} = \qty{3.719 \pm 0.001}{ppm}$.

\textbf{16.} Het kwadratisch gemiddelde residu is \qty{0.0014}{ppm}, de
grootte van de afleesprecisie, zonder trend: het 1:1-model past.

\textbf{17.} Gebonden fractie $f$ bij $G_0 = fH_0 + f/(K(1 - f))$: 0.28 equivalent voor $f = 0.2$ en 2.1 voor $f = 0.8$.

\textbf{18.} $KH_0 = 10^{6.1} \times \num{2.0e-3} \approx 2500$: de kromme zou scherp knikken bij één equivalent
en alleen een ondergrens geven; er is een competitie-experiment of
calorimetrie bij lagere concentratie nodig.

\textbf{19.} In het tolueen, een homogene oplossing van alkeen en
permanganaat.

\textbf{20.} Het gereduceerde mangaan verlaat de organische fase (als
mangaandioxide), en de \omterm{def:b3:supramolecular:hosts}{kroonether}, vrijgekomen aan het grensvlak, neemt
een nieuw $\ce{K+}$ en $\ce{MnO4-}$ op. Elk kroonethermolecuul
vervoert vele keren permanganaat: een katalytische hoeveelheid volstaat.

\textbf{21.} $K_{\mathrm{ex}}$ vermenigvuldigt $[\ce{K+}]$: een hoge kaliumconcentratie duwt de extractie vooruit.

\textbf{22.} Quaternaire ammoniumzouten, zoals tetrabutylammonium- of
benzyltriethylammoniumchloride, waarvan het kation op zich lipofiel is.

\textbf{23.} Benzeen veroorzaakt leukemie; tolueen doet hetzelfde werk met
een veel kleiner gevaar.

\textbf{24.} Met \qty{1.0}{mmol/L} \omterm{def:b3:supramolecular:hosts}{kroonether} wordt 73\,\% van het
permanganaat in het tolueen geëxtraheerd.
\end{solution}
